
making pandemic-related data computationally tractable.

The history of humankind has been marked by deadly pathogen-borne epidemics which has contributed to radically shifting our species today, from our genes to our behaviour. In our genes, viral DNA, immune system etc. In our behaviour, learned hygiene etc. A major cause of epidemics is human-animal interaction, which have been catalyzed by the advent of domestication, as well as technological advancements which have brought us in closer contact with unfriendly animals. Although tremendous progress in science, medicine, as well as public health awareness have helped curtail the onset and the impact of epidemics, the acceleration of globalisation and climate change in the 21st century are repeatedly pointed out as a trigger for an increased frequency of global-scale pandemics~\cite{marani2021}. ecosystem modifications (e.g., deforestation, loss of biodiversity~\cite{keesing2010}) reduce the distances between zoonotic reservoirs and human populations, and disease-carrying vectors (e.g., mosquitoes~\cite{jones2008, ryan2019} and rodents~\cite{johnson2020}) are able to adapt and proliferate in ever-larger regions of the globe due to global warming. Following the outbreak of \gls{sars} in the early 2000s, non-governmental global organisations such as the \gls{who} have introduced measures to strengthen international coordination to better the management of pandemics. In addition, scientific community has called for a transition from a treatment-led to a prevention-led health system, and increasing research efforts for \emph{pandemic preparedness}. The recent COVID-19 pandemic has illustrated that preventing and preparing against pandemics is a Herculean task. Beyond `para-scientific` issues such as international coordination, supply chain management and public health education, it is an incredibly multi-faceted problem from a scientific perspective~\cite{williams2023}. Microbiology for microbes, ecology for animals, evolutionary biology to understand pathogen evolution, network science to understand transmission, algorithmics, machine learning, and artificial intelligence for modelling and efficient computing~\cite{kraemer2025}. In the past decades, data science has emerged as a cornerstone of pandemic preparedness, as preparedness measures are driven by the collection of analysis of large-scale datasets, such as wildlife monitoring footage, pathogenic genomes from wastewater, human mobility patterns, or electronic health records~\cite{grenfell2004, khurana2024}. Thus, investment in data-driven research is crucial for our preparedness against the best pandemic. I divide the research perspectives as `exploratory`, i.e., analysing the potential of new forms of data or analysis methods, or `exploitatory`, i.e., optimizing existing aspects of the pandemic data analysis cycle. This could encompass the development of aggregating platforms for data or methods sharing, data and software engineering, or methods for decentralized analyses. Dividing

This thesis can be viewed as a data science's perspective on pandemic preparedness. In particular, my thesis focuses on making pandemic-related data more computer-friendly \textcolor{red}{vector?}. In the `exploitatory` domain, I introduce phylo2vec~\cite{penn2025phylo2vec, scheidwasser2025}, a method for encoding phylogenetic trees (\textcolor{red}{definition}) into a simple integer-based vector (e.g., a list of numbers). The method can be viewed as a data engineering method, compressing data from phylogenetic trees into simple objects. In the `exploratory` domain, I analyse the potential of two forms of data, protein structures and video-based tracking, to respectively better understand the evolution of pathogens and detect infection early on and non-invasively.

%, which have favoured the spillover of zoonoses~\cite{sikora2025}. These 

Human populations have always faced deadly epidemics throughout history. A major cause for such episodes are human-animal interactions, which have favoured the spillover of zoonoses~\cite{sikora2025}. Learned progress in medicine (black plague, john snow), but the 21st century is marked by an acceleration of global warming and climate change. Amongst the many challenges that humanity will have to face, environmental experts expect an increased risk of pandemics. There are multiple factors behind this: ecosystem changes such as deforestation for crops, mining, and roads reduce the distances between wild animal reservoirs and human populations - disease-carrying vectors have more territory (e.g., rodents~\cite{johnson2020}, mosquitoes~\cite{jones2008, ryan2019}), globalisation favours transmission. Loss of biodiversity~\cite{keesing2010}. After the SARS outbreak, non-governmental global organisations such as the \gls{who} and the World Bank have set up teams/divisions/reports revolving around pandemics. Mention Disease X. The scientific community has called for a transition from a treatment-led to a prevention-led health system. 

\begin{itemize}
\item SARS-CoV-2 pandemics + H5N1 pandemic show the many degrees of freedom that surround pandemic preparedness.
\item In addition to international coordination, logistics, preventive measures, there are many scientific aspects to studying epidemics/pandemics. 
\item SARS-CoV-2 has for instance shown the importance of large-scale sampling of viral sequences (\cite{khurana2024}) to study the virus: its evolution, its transmissibility to humans etc.
\item Analysing these heaps of sequences requires sophisticated bioinformatics pipelines, efficient software, and good hardware. To understand evolutionary relationships, phylogenetic methods have become more and more efficient. The object, a phylogenetic tree, is the common object to understand relationships between entities. However, inferring phylogenetic trees involves many subcomponents which are not-trivial: how to sample a random tree, how to represent this tree for a computer (and a human) to understand it etc. While deep learning has been a revolution for vision ang language, it has not been successful yet for the most important analysis components of phylogenetics \cite{sapoval2022}.
\item Sequences alone are only part of the problem: biomechanics are also important to understand how the pathogen infects its host/target, e.g., binding with ACE for SARS-CoV-2 and humans. crucial for vaccine design. This requires advanced experimental techniques such as x-ray crystallography, Cryo-EM, or protein sturcutre prediction models trained on millions of proteins (AlphaFold \cite{jumper2021}, ESM \cite{lin2023}). Whereas sequence-data is the default for understanding the evolution of pathogens via traditional phylogenetic analyses, protein structures is the thing that actually is there in real life.
\item While sequencing has become increasingly cheap and fast over the past decades, this is not scalable for the surveillance of wildlife or livestock, which are a huge contributor to zoonotic epidemics (which are the majority of epidemics). This is especially the case for small animals such as birds, which have contributed to avian influenza epidemics, or bats (coronaviruses). A potential solution in that consists in monitoring and analysing the behaviour of animal reservoirs. This is done to some extent with humans who go to fields and farms, but a potential cool avenue would consist in video surveillance/monitoring. Being able to collect large-scale datasets of those is crucial for collaboration, and easily scalable without needing much human intervention. However, analysing such big quanities of data is difficult, so using computer vision techniques might help assist humans in this task. Beyond traditional computer vision techniques such as optical flow, which has been successfully implemented in broiler chicken farms to measure onset of Campylobacter epidemics~\cite{colles2016}).
\item why vectorized? Because of the new image of informatics: machine lerning, deep learning, parallel computing, very fast GPUs. GPUs drink vectors for breakfast, and vectors are among the most efficient representations of data. Moore's law. 
\item My contributions are as follows: a vector-based representation of phylogenetic trees, called Phylo2Vec, a phylogenetic analysis of the evolution of spike proteins in coroanviruses using protein structures using a vector-quantized encoding of tertiary structure, and a pilot system for early detection of infection in chicken using deep learning-based video tracking
\end{itemize}